\scanpage{10}
% L1–4
$X,Y$ deux espaces topologiques, considérons $\mathrm{Top}(X)\mathbin{\underset{\mathrm{top}}{\overset{2}{\times}}}\mathrm{Top}(Y)$ et $\mathrm{Top}(X\times Y)$ (on vérifie facilement que ce dernier ne dépend que de $\mathrm{Top}(X)$ et de $\mathrm{Top}(Y)$, i.e. des \deleted{\unsure{les}} espaces sobres associés à $X$ et $Y$) On a un morphisme canonique de topos

% L5
\[
\deleted{\unsure{\mathrm{Top}|\mathrm{Prop.}}}\quad\tag*{\ensuremath{(*)}}\boxed{\mathrm{Top}(X\times Y)\longrightarrow\mathrm{Top}(X)\mathbin{\underset{\mathrm{top}}{\overset{2}{\times}}}\mathrm{Top}(Y)}
\]

% L6
et on veut des conditions pour que ce soit une équivalence.

% L6-left
\margin{$O_X=$ cat. des ouverts de $X$}

% L7–14
Or on voit que $\mathrm{Top}(X\times Y)=\widetilde{O_X\times O_Y}$, où la topologie à mettre sur $O_X\times O_Y$ est la ``topologie produit'' \inserted{$\pi$}, i.e. celle pour laquelle les faisceaux sont les foncteurs contravariants $(O_X\mathbin{\overset{\deleted{\unsure{o|2}}}{\times}}O_Y)^o\longrightarrow(\mathrm{Ens})$ qui sont des faisceaux \uline{séparément}. \deleted{Il} Il y a aussi la topologie $\pi'$ \deleted{\unsure{par}} ``induite'' par le produit $X\times Y$ en utilisant l'injection (supposant $X$ \deleted{\unsure{et}} $Y$ \deleted{\unsure{=}} $\ne\varnothing$) $O_X\times O_Y\subset O_{X\times Y}$. Il est clair que tout faisceau pour cette topologie (faisceau sur $X\times Y$) est un bifaisceau \deleted{\unsure{pour donc}} donc la top. $\pi'$ est plus fine que $\pi$.

% L12-left
\margin{\uline{NB} les $U\times V$ ($U\in O_X$, $V\in O_Y$) font une \uline{base} de la top. de $X\times Y$}

% L14–19
On voit déjà que le \deleted{\unsure{topos}} morphisme $(*)$ est un \uline{plongement de topos}, et on est ramené à voir si ce plongement est une équivalence i.e. $\pi=\pi'$. Donc on est amené à voir si tout recouvrement pour $\pi'$ de $U\times V$ ($U\in O_X$, $V\in O_Y$) par des $U_i\times V_i$ [i.e. rec. ensembliste dans $X\times Y$) est un rec. pour $\pi$, i.e. \deleted{si} si

% L15–end-left
% slanted
\margin{\uline{NB} Comme $\mathrm{Top}(X\times Y)$ est un sous-topos de $\mathrm{Top}(X)\times\mathrm{Top}(Y)$, ayant \uline{suffisamm}\unsure{t} de pts, et contenant tous les pts de $\mathrm{Top}(X\times Y)$, on en conclut que c'est le plus \unsure{petit} sous topos de $\mathrm{Top}(X)\times\mathrm{Top}(Y)$ contenant les pts du dit topos. Analogue pour produit d'une famille quelconque. Donc la question si \unsure{$(*)$|c'}est une équivalence \uline{équivaut} à celle si $\mathrm{Top}(X)\mathbin{\underset{\mathrm{top}}{\overset{2}{\times}}}\mathrm{Top}(Y)$ a assez de points\ldots}

% L20
\[
\bigcup_i U_i\times V_i=U\times V\quad\Longrightarrow\quad\mathop{\mathrm{Sup}}_i U_i\times V_i=U\times V
\]

% L21–22
le sup étant pris dans l'ens des sous objets de $U\times V$ dans le topos $\widetilde{O_X\times O_Y}$ (top $\pi$).

% L23–25
\uline{Proposition} \quad Pour qu'il en soit ainsi, il suffit que tout point de $U$ (\uline{ou} de $V$) admette un voisinage (pas néc. ouvert) quasi-compact

% L26–27
\uline{Corollaire} \quad \deleted{\unsure{Si}} \deleted{\unsure{tout|tous}} pt de $X$ \inserted{(ou de $Y$)} admet un syst. fond. de voisinages qui sont quasi-compacts, alors $(*)$ est une équivalence.

% foot
% archive sheet number: 9
